% ============================
% Classe du document
% ============================
\documentclass[ % 'report' est bien adapté aux thèses, 'memoir' plus complet (equivalents 'scrbook', 'scrreprt' de Koma-script))
a4paper,
11pt,
twoside
]{memoir}

% ============================
% Chargement des réglages
% ============================
% --- Langue et encodage (gérés par fontspec si LuaLaTeX/XeLaTeX) ---
\usepackage[english, french]{babel}         % Gestion multilingue, langue française par défaut (dernière) : typographie, noms automatiques, utiliser \selectlanguage{english} pour basculer en anglais et \begin{otherlanguage}{french} pour des inclusions fr suivi de \end{otherlanguage} puis de \selectlanguage{english} pour reprendre une partie en anglais.
\usepackage{csquotes}                       % Guillemets typographiques adaptés à la langue
\addto\captionsfrench{%
  \renewcommand{\listfigurename}{Liste des illustrations}%
  \renewcommand{\listtablename}{Liste des tableaux}%
}
\usepackage{multicol}                       % pour un environnement 2-colonnes
% A supprimer lorsque j'utiliserai des polices personnalisées.
\usepackage[T1]{fontenc}                    % Sortie correcte des accents
\usepackage{lmodern}                        % Police vectorielle moderne

% --- Mathématiques et symboles scientifiques ---
\usepackage{amsmath, amssymb, amsfonts}     % Pour les formules et symboles mathématiques
\usepackage{siunitx}                        % Gestion des unités scientifiques (1.23 en et 1,23 fr) \SI{3.5}{mg/L} -> possibilité aussi de faire des unités personnalisées voir chatGPT au besoin

% --- Images et figures ---
\usepackage{graphicx}                       % Insertion d'images
\usepackage{float}                          % Placement forcé avec [H]
\usepackage[font=small,labelfont=bf,width=0.8\textwidth]{caption}   % Style des légendes
\usepackage{subcaption}                     % Sous-figures
% \usepackage{svg}                          % Pour inclure des fichiers SVG

% --- Tableaux et couleurs ---
\usepackage[table,xcdraw]{xcolor}           % Couleurs dans les tableaux et texte
\usepackage{multirow}                       % Fusion de cellules
% \usepackage{booktabs}                     % Tableaux élégants si besoin (\toprule...)

% --- Bibliographie ---
\usepackage[backend=biber, style=apa, defernumbers=true, refsection=chapter, sorting=nyt, url=false]{biblatex} % style=iso-alphabetic, cas complexes ! norme AFNOR Z 44005 (reco nationale)
\addbibresource{references.bib}             % Fichier de bibliographie
% \addbibresource[label=art1]{biblio_article1.bib}      % Autre fichier possible -> utiliser betterbiblatex
% Réglages typographiques utiles
\setcounter{biburllcpenalty}{7000}
\setcounter{biburlucpenalty}{8000}
%Augmente le seuil de coupure dans les URL en bibliographie. Cela évite les césures moches dans les longues URL.
%\renewcommand*{\labelnamepunct}{\newunitpunct\par}% Met un saut de ligne après les noms d’auteur au lieu d’un point, pour une bibliographie plus aérée.
%\renewbibmacro{in:}{\newline In:}           % Met également un saut de ligne avant le mot In: (chapitre dans un ouvrage), pour un rendu plus clair.
%\DefineBibliographyExtras{french}{\restorecommand\mkbibnamefamily} % Restaure un comportement cohérent en français pour la bibliographie dans la mise en forme des noms de famille

% --- Glossaire et acronymes ---
%\usepackage[toc, acronym]{glossaries}
\usepackage[acronym, toc]{glossaries-extra}
\glssetcategoryattribute{acronym}{counter}{none} % désactive les compteurs
\setabbreviationstyle[acronym]{long-short}
\glsdisablehyper % <-- ici : désactivation des hyperliens
\makeglossaries
% ---------- MOTS ---------- %
\newglossaryentry{accession fine}{
    name={Accession fine}, 
    text={accession fine}, 
    description={Accession ou lignée de vigne donnant des grappes millerandées, c'est-à-dire avec des baies de tailles variables et un port lâche. Plus le degré de finesse est élevé, plus l'hétérogénéité est grande},
    plural={accessions fines}
}

\newglossaryentry{baguette}{
    name={Baguette}, 
    text={baguette}, 
    description={Terme viticole, notamment dans la taille en Guyot, désignant le bois de l'année précédente laissé sur le pied pour porter les rameaux fructifères de l'année suivante},
    plural={baguettes}
}

% ---------- ACRONYMES ---------- %
\newacronym{aba}{ABA}{Acide abscissique}
\newacronym{acp}{ACP}{Analyse en composantes principales}

% --- Commentaires en marge (optionnels) ---
\usepackage[textsize=footnotesize, color=orange!20]{todonotes}

% --- Mini ToC pour chaque chapitres ---
%\usepackage{minitoc} % Compatible avec memoir, sinon etoc, sinon fonction interne (voir manuel p.415)

% --- Hyperliens ---
\usepackage[colorlinks=true,
            linkcolor=black,
            citecolor=black,
            filecolor=black,
            urlcolor=blue]{hyperref}
%\usepackage[hidelinks]{hyperref}            
% \usepackage{cleveref}                     % Références intelligentes (ex: \cref{fig:truc})               % chargement des packages utilisés
%% Utilisation de polices système/personnalisées avec fontspec
% Pour utiliser fontspec (nécessite xelatex ou lualatex) -> Xualatex mieux
\usepackage{fontspec}

% Police principale (corps du texte)
\setmainfont{RobotoSerif}[% Le nom exact peut dépendre du système, vérifiez les fichiers .ttf/.otf
  Path=Fonts/Roboto_Serif/static/,
  UprightFont={*_72pt-Light},    % Pour le texte normal (Light 300)
  ItalicFont={*-LightItalic}, % Pour l'italique (Light 300 Italic)
  BoldFont={*-SemiBold},        % Pour le gras (Regular)
  BoldItalicFont={*-SemiBoldItalic},    % Pour le gras italique (Italic)
  Extension=.ttf
]

% Police pour les titres : Roboto Slab Light 300
\newfontfamily\titrefont[
    Path=Fonts/Roboto_Slab/static/,
    UprightFont = *-Light.ttf,
    BoldFont = *-Regular.ttf
]{RobotoSlab}

% Couleurs globales du document : éclaircir les polices au besoin
%\definecolor{darkgraytext}{gray}{0.15}
%\color{darkgraytext}

% Polices des titres de section (Roboto Slab Light 300)
% memoir utilise des commandes spécifiques pour les polices des titres de section
%\renewcommand{\chapnamefont}{\normalfont\fontspec{Roboto Slab}[UprightFont={*-Light300}]} % Nom du chapitre (si affiché)
%\renewcommand{\chapnumfont}{\normalfont\fontspec{Roboto Slab}[UprightFont={*-Light300}]} % Numéro de chapitre (si affiché)
%\renewcommand{\chaptitlefont}{\normalfont\fontspec{Roboto Slab}[UprightFont={*-Light300}]\Huge} % Titre de chapitre
%\renewcommand{\sectfont}{\normalfont\fontspec{Roboto Slab}[UprightFont={*-Light300}]\Large} % Titre de section
%\renewcommand{\subsectfont}{\normalfont\fontspec{Roboto Slab}[UprightFont={*-Light300}]\large} % Titre de sous-section
%\renewcommand{\subsubsectfont}{\normalfont\fontspec{Roboto Slab}[UprightFont={*-Light300}]\normalsize} % Titre de sous-sous-section
% Et ainsi de suite pour les autres niveaux si vous les utilisez (paragraph, subparagraph)
%\renewcommand{\chaptitlefont}{\titlefont\LARGE} % Titre
% tailles polices : \tiny < \scriptsize < \footnotesize < \small < \normalsize < \large < \Large < \LARGE < \huge < \Huge
% graisse : \bfseries

% Ne pas oublier de recaluler la largeur des blocs texte après un changement de police, de graisse de police ou même de taille !
% Bien tout définir...                  % utilisation de polices personnalisées
\usepackage[pass]{geometry} % Pour des mises en forme local dans mon doc sans prise de tête liminaires et articles
% ---- Utilisation des fonctionnalités de la classe memoir ----
% 1. Taille du papier (définie lors de l'appel à la class memoir
%\settrims{0pt}{0pt} %\settrims{<haut>}{<bord_de_coupe>} en cas de bords perdus, en définir 2

% 2. Bloc de texte & marges
%\settypeblocksize{<hauteur>}{<largeur>}{<ratio>}
%\setlrmargins{⟨spine⟩}{⟨edge⟩}{⟨ratio⟩}
%\setulmargins{<haut>}{<bas>}{<ratio>} % Utilisez cette commande si la hauteur de votre bloc de texte est déjà fixe (par \settypeblocksize) et que vous positionnez le bloc en ajustant les marges, ex : \settypeblocksize{41\baselineskip + \topskip}{33pc}{*}
\setlrmarginsandblock{2.5cm}{5cm}{*} % Marges latérales (gauche/droite)
\setulmarginsandblock{3cm}{3cm}{*} % Marges haut/bas

% 3. En-têtes et pieds de page
\setheadfoot{\onelineskip}{2\onelineskip} % \setheadfoot{<hauteur_entete>}{<espacement_pied_de_page>}
%\setheaderspaces{1.5\onelineskip}{*} % \setheaderspaces{<chute_entete>}{<separation_entete_texte>}{<ratio>}
\setlength{\headheight}{15pt}

% 4. Notes marginales
\setmarginnotes{1em}{3.5cm}{1.5em} %\setmarginnotes{<separation>}{<largeur>}{<poussee>}

% Interligne plus aéré
\linespread{1.05}   % Memoir intègre en interne la gestion de l'interlignage...

% Paragraphe : pas d'indentation + saut entre paragraphes
\setlength{\parindent}{1em} % inutile
\setlength{\parskip}{0.2\onelineskip} %\setlength{\parskip}{0.7em}

% Fin. Corrections automatiques de la mise en page
%\checkandfixthelayout[nearest] % Ajuste la hauteur du bloc de texte à un multiple entier de lignes
\checkandfixthelayout % application immédiate

% Style de page minimaliste personnalisé
\makepagestyle{minimaliste}
\makeevenhead{minimaliste}{}{}{\itshape\leftmark} % chapitre à droite (pages paires)
\makeoddhead{minimaliste}{\itshape\rightmark}{}{} % section à gauche (pages impaires)
\makeevenfoot{minimaliste}{\thepage}{}{} % numéro page à gauche en bas (pages paires)
\makeoddfoot{minimaliste}{}{}{\thepage} % numéro page à droite en bas (pages impaires)

% Pages paires : numéro en bas à gauche
\makeevenfoot{plain}{\thepage}{}{}
% Pages impaires : numéro en bas à droite
\makeoddfoot{plain}{}{}{\thepage}

% Style des titres (sobre) % maîtrise fine
%\usepackage{titlesec}
%\titleformat{\chapter}[hang]{\Large\bfseries\titlefont}{}{0pt}{}
%\titleformat{\section}[hang]{\large\bfseries}{}{0pt}{}
%\titleformat{\subsection}[runin]{\bfseries}{}{0pt}{}
%\titlespacing*{\chapter}{0pt}{20pt}{10pt}
%\titlespacing*{\section}{0pt}{12pt}{6pt}

% Style des titres natif à memoir (simple et stable)
\chapterstyle{veelo} % ou 'default', 'veelo', 'companion', 'ell', 'demo2', 'dash', etc.

\addto\captionsfrench{%
  \renewcommand{\listfigurename}{Liste des illustrations}%
  \renewcommand{\listtablename}{Liste des tableaux}%
}

% Mettre paragraph et subparagraph sur leur propre ligne
%\setbeforeparaskip{1.5ex}   % espace avant paragraph
%\setafterparaskip{1.0ex}    % espace après paragraph
%\setparaheadstyle{\normalfont\bfseries}    % paragraphe en gras
%\setbeforeparheadhook{}   
%\setafterparheadhook{\par}  % forcer saut de ligne

%\setbeforesubparaskip{1.2ex}
%\setaftersubparaskip{0.8ex}
%\setsubparaheadstyle{\normalfont\itshape}
%\setbeforesubparaheadhook{} 
%\setaftersubparaheadhook{\par}

% ---- Cadre encart remerciements ----
\tcbset{
  gratitude/.style={
    colframe=gray!50,        % couleur du contour
    colback=white,           % fond
    boxrule=0.3pt,           % épaisseur du trait
    arc=2pt,                 % coins légèrement arrondis
    left=8pt, right=8pt, top=6pt, bottom=6pt,
    width=0.75\textwidth,    % largeur fixe
    enhanced jigsaw,         % rendu propre
    before skip=1cm,         % espace au-dessus
    after skip=0pt,          % espace en-dessous
    overlay={%
      \node[anchor=north west, xshift=-0.6cm, yshift=0.6cm] at (frame.north east)
        {\includegraphics[width=0.9cm]{0_Liminaires/Images/EtoilePP.png}}; % étoile qui déborde
    }
  }
}                 % mise en page
% Pour les commandes personnalisées

% --- Page blanche entre sections, sans numéro ---
\newcommand{\blankpage}{%
  \newpage
  \null
  \thispagestyle{empty}%
  \addtocounter{page}{-1}%
  \newpage
}

% --- Saut vertical court (optionnel) ---
\newcommand{\smallskipline}{\vspace{0.5\baselineskip}}

% --- Saut vertical moyen (optionnel) ---
\newcommand{\medskipline}{\vspace{1\baselineskip}}

% --- Saut vertical long (optionnel) ---
\newcommand{\bigskipline}{\vspace{2\baselineskip}}

% --- Ligne horizontale (utile en page de garde, etc.) ---
\newcommand{\HRule}{\rule{\linewidth}{0.4pt}}

% --- Commentaire en marge (léger) ---
\usepackage[textsize=footnotesize, color=orange!20]{todonotes}
\newcommand{\commentaire}[2][]{\todo[#1]{#2}}

% --- Paragraphe mis en valeur (style "new thought") ---
\newcommand{\newthought}[1]{\textbf{#1}\hspace{0.25em}}

% --- Environnement pour forcer une page droite non paginée ---
\newenvironment{rightpage}{%
  \cleardoublepage      % va à la page droite suivante
  \thispagestyle{empty} % pas de numéro
}{%
  \vfill                % pousse le contenu en haut
  \cleardoublepage      % termine sur une page droite vierge
}               % commandes personnalisées

\begin{document}

% Page de garde et liminaires
\pagestyle{empty}

% Voir recommandations ED SVS !!!!!! puis adapter

\pagestyle{empty}

% Les logos en tête

\includegraphics[scale=1, height=1.7cm]{0_Liminaires/Images/UBX_logo.png}
\hfill

\includegraphics[scale=1, height=1.7cm]{0_Liminaires/Images/RF.png} %  remplacer par logo UMR (+ ISVV ?)
\hfill

\begin{center}
\doublespacing
\begin{Large}

THÈSE PRÉSENTÉE\\ POUR OBTENIR LE GRADE DE \\
{\LARGE \textbf{DOCTEUR\\DE L'UNIVERSITÉ DE BORDEAUX} } \\
\vspace{0.55cm}
ECOLE DOCTORALE SCIENCES DE LA VIE ET DE LA SANTÉ\\
{\normalsize Spécialité OENOLOGIE} \\
\vspace{0.55cm}
Par \textbf{Guillaume Séraphin BLANC LELEU} \\
\vspace{0.55cm}
{\Large Titre.}
\end{Large}
\vspace{0.55cm}
\begin{normalsize}
\begin{singlespace}
Sous la direction de : \textbf{Tristan RICHARD}\\
Responsable scientifique : \textbf{...}
\end{singlespace}
\end{normalsize}
\end{center}
\vfill
\begin{center}
{\large Soutenue le x novembre 2025 }\\    
\end{center}
\vfill
Membres du jury :
\begin{table}[b]
\centering
\makebox[\textwidth]{%
\begin{tabular}{lllr}
Mme Prénom NOM & Grade & Université & Rapporteuse \\
Mme Prénom NOM & Grade & Université & Fonction \\
Mme Prénom NOM & Grade & Université & Fonction \\
Mme Prénom NOM & Grade & Université & Fonction \\
Mme Prénom NOM & Grade & Université & Fonction \\
\end{tabular}
}
\end{table}

\thispagestyle{empty}
\vspace*{0pt}
\vfill

\afterpage{\blankpage}
\newpage

\begin{small}
\begin{singlespace}
\begin{center}
\textbf{Titre.}
\end{center}
\textbf{Résumé :} 
...

L.\newline
\\
\textbf{Mots-clés :} ... \\

%\noindent\rule{\textwidth}{1pt}
\noindent\makebox[\linewidth]{\rule{\textwidth}{0.4pt}}

\vfill
% \selectlanguage{french}
\begin{center}
    \textbf{Unité de recherche}\\
UMR 1366 OENO, 33140 Villenave-d'Ornon, France.
\end{center}

\afterpage{\blankpage}
\newpage
% \selectlanguage{english}

\begin{center}
\textbf{Title english.}
\end{center}
\textbf{Abstract:} 
...
\\
\newline
\textbf{Keywords:} ... \\

\noindent\makebox[\linewidth]{\rule{\textwidth}{0.4pt}}

\vfill
% \selectlanguage{french}
\begin{center}
    \textbf{Reasearch unit}\\
UMR 1366 OENO, 33140 Villenave-d'Ornon, France.
\end{center}
\end{singlespace}
\end{small}
\frontmatter
\addcontentsline{toc}{section}{Citations}
\section*{Citations}
\textit{"Blablabla citation 1"}
\begin{flushright}
Auteur 1.
\end{flushright}

\vspace{10pt}

\textit{"Blablabla citation 2"}
\begin{flushright}
Auteur 2.
\end{flushright}


\addcontentsline{toc}{section}{Remerciements}
\section*{Remerciements}

Fait toi plaisir quand tu en seras là c'est que tu auras déjà bien avancé ;)

% Maintenant déballe ton CV
\newpage
\addcontentsline{toc}{section}{Productions scientifiques}
\section*{Productions scientifiques}
\subsection*{Communications orales}
\begin{itemize}
    \item Workshop Modélisation 2024 - Toulouse (17/03/2024)
    \item ...
\end{itemize}

\subsection*{Posters scientifiques}

\begin{itemize}
    \item Journées Jeunes Chercheurs 2023 - Clermont-Ferrand (03/04/2023)
    \item ...
\end{itemize}

\vspace{10pt}

\subsection*{Activités annexes à la recherche}
\begin{itemize}
    \item Encadrement de stagiaires
    \item ADPI...
\end{itemize}

% Table des matières
\pagestyle{plain}
\tableofcontents

% Glossaire (généré si nécessaire)
% Mise en forme du glossaire contenu dans glossaire_definitions.tex

%\chapter*{Glossaire} % crée un chapitre non numéroté, il n’est pas ajouté automatiquement à la table des matières
%\addcontentsline{toc}{chapter}{Glossaire} % Cette commande ajoute manuellement une ligne dans la table des matières (ToC). toc = table of contents ; chapter = niveau hiérarchique de l’entrée (ici, un chapitre) ; Glossaire = le texte à afficher dans la ToC

% Glossaire et acronymes — à inclure avec % Mise en forme du glossaire contenu dans glossaire_definitions.tex

%\chapter*{Glossaire} % crée un chapitre non numéroté, il n’est pas ajouté automatiquement à la table des matières
%\addcontentsline{toc}{chapter}{Glossaire} % Cette commande ajoute manuellement une ligne dans la table des matières (ToC). toc = table of contents ; chapter = niveau hiérarchique de l’entrée (ici, un chapitre) ; Glossaire = le texte à afficher dans la ToC

% Glossaire et acronymes — à inclure avec % Mise en forme du glossaire contenu dans glossaire_definitions.tex

%\chapter*{Glossaire} % crée un chapitre non numéroté, il n’est pas ajouté automatiquement à la table des matières
%\addcontentsline{toc}{chapter}{Glossaire} % Cette commande ajoute manuellement une ligne dans la table des matières (ToC). toc = table of contents ; chapter = niveau hiérarchique de l’entrée (ici, un chapitre) ; Glossaire = le texte à afficher dans la ToC

% Glossaire et acronymes — à inclure avec \include{Liminaires/glossaire}

% Affiche la liste des acronymes
\printglossary[
  type=\acronymtype,
  title=Acronymes,
  toctitle=Acronymes % Alternative : Liste des acronymes, changera le titre uniquement dans la toc (table of content)
]

% Affiche le glossaire principal
\printglossary[
  title=Glossaire,
  toctitle=Glossaire % Glossaire des termes
]

% Utilisation [à supprimer]
Une fois un mot ou acronyme déclaré avec \newglossaryentry ou \newacronym :
Pour l’utiliser dans ton texte :
\gls{nmr} → affiche NMR + explication complète la 1re fois (si configuré)

\Gls{nmr} → idem, mais avec majuscule

\glspl{nmr} → pluriel

\acrshort{nmr} → affiche juste NMR

\acrfull{nmr} → affiche acide abscissique (ABA) (forme longue + sigle)

\acrlong{nmr} → juste la forme longue

LaTeX se charge automatiquement d’écrire le mot complet la première fois, puis seulement l’abréviation ensuite.

% Affiche la liste des acronymes
\printglossary[
  type=\acronymtype,
  title=Acronymes,
  toctitle=Acronymes % Alternative : Liste des acronymes, changera le titre uniquement dans la toc (table of content)
]

% Affiche le glossaire principal
\printglossary[
  title=Glossaire,
  toctitle=Glossaire % Glossaire des termes
]

% Utilisation [à supprimer]
Une fois un mot ou acronyme déclaré avec \newglossaryentry ou \newacronym :
Pour l’utiliser dans ton texte :
\gls{nmr} → affiche NMR + explication complète la 1re fois (si configuré)

\Gls{nmr} → idem, mais avec majuscule

\glspl{nmr} → pluriel

\acrshort{nmr} → affiche juste NMR

\acrfull{nmr} → affiche acide abscissique (ABA) (forme longue + sigle)

\acrlong{nmr} → juste la forme longue

LaTeX se charge automatiquement d’écrire le mot complet la première fois, puis seulement l’abréviation ensuite.

% Affiche la liste des acronymes
\printglossary[
  type=\acronymtype,
  title=Acronymes,
  toctitle=Acronymes % Alternative : Liste des acronymes, changera le titre uniquement dans la toc (table of content)
]

% Affiche le glossaire principal
\printglossary[
  title=Glossaire,
  toctitle=Glossaire % Glossaire des termes
]

% Utilisation [à supprimer]
Une fois un mot ou acronyme déclaré avec \newglossaryentry ou \newacronym :
Pour l’utiliser dans ton texte :
\gls{nmr} → affiche NMR + explication complète la 1re fois (si configuré)

\Gls{nmr} → idem, mais avec majuscule

\glspl{nmr} → pluriel

\acrshort{nmr} → affiche juste NMR

\acrfull{nmr} → affiche acide abscissique (ABA) (forme longue + sigle)

\acrlong{nmr} → juste la forme longue

LaTeX se charge automatiquement d’écrire le mot complet la première fois, puis seulement l’abréviation ensuite.
%% Mise en forme du glossaire contenu dans glossaire_definitions.tex

%\chapter*{Glossaire} % crée un chapitre non numéroté, il n’est pas ajouté automatiquement à la table des matières
%\addcontentsline{toc}{chapter}{Glossaire} % Cette commande ajoute manuellement une ligne dans la table des matières (ToC). toc = table of contents ; chapter = niveau hiérarchique de l’entrée (ici, un chapitre) ; Glossaire = le texte à afficher dans la ToC

% Glossaire et acronymes — à inclure avec % Mise en forme du glossaire contenu dans glossaire_definitions.tex

%\chapter*{Glossaire} % crée un chapitre non numéroté, il n’est pas ajouté automatiquement à la table des matières
%\addcontentsline{toc}{chapter}{Glossaire} % Cette commande ajoute manuellement une ligne dans la table des matières (ToC). toc = table of contents ; chapter = niveau hiérarchique de l’entrée (ici, un chapitre) ; Glossaire = le texte à afficher dans la ToC

% Glossaire et acronymes — à inclure avec % Mise en forme du glossaire contenu dans glossaire_definitions.tex

%\chapter*{Glossaire} % crée un chapitre non numéroté, il n’est pas ajouté automatiquement à la table des matières
%\addcontentsline{toc}{chapter}{Glossaire} % Cette commande ajoute manuellement une ligne dans la table des matières (ToC). toc = table of contents ; chapter = niveau hiérarchique de l’entrée (ici, un chapitre) ; Glossaire = le texte à afficher dans la ToC

% Glossaire et acronymes — à inclure avec \include{Liminaires/glossaire}

% Affiche la liste des acronymes
\printglossary[
  type=\acronymtype,
  title=Acronymes,
  toctitle=Acronymes % Alternative : Liste des acronymes, changera le titre uniquement dans la toc (table of content)
]

% Affiche le glossaire principal
\printglossary[
  title=Glossaire,
  toctitle=Glossaire % Glossaire des termes
]

% Utilisation [à supprimer]
Une fois un mot ou acronyme déclaré avec \newglossaryentry ou \newacronym :
Pour l’utiliser dans ton texte :
\gls{nmr} → affiche NMR + explication complète la 1re fois (si configuré)

\Gls{nmr} → idem, mais avec majuscule

\glspl{nmr} → pluriel

\acrshort{nmr} → affiche juste NMR

\acrfull{nmr} → affiche acide abscissique (ABA) (forme longue + sigle)

\acrlong{nmr} → juste la forme longue

LaTeX se charge automatiquement d’écrire le mot complet la première fois, puis seulement l’abréviation ensuite.

% Affiche la liste des acronymes
\printglossary[
  type=\acronymtype,
  title=Acronymes,
  toctitle=Acronymes % Alternative : Liste des acronymes, changera le titre uniquement dans la toc (table of content)
]

% Affiche le glossaire principal
\printglossary[
  title=Glossaire,
  toctitle=Glossaire % Glossaire des termes
]

% Utilisation [à supprimer]
Une fois un mot ou acronyme déclaré avec \newglossaryentry ou \newacronym :
Pour l’utiliser dans ton texte :
\gls{nmr} → affiche NMR + explication complète la 1re fois (si configuré)

\Gls{nmr} → idem, mais avec majuscule

\glspl{nmr} → pluriel

\acrshort{nmr} → affiche juste NMR

\acrfull{nmr} → affiche acide abscissique (ABA) (forme longue + sigle)

\acrlong{nmr} → juste la forme longue

LaTeX se charge automatiquement d’écrire le mot complet la première fois, puis seulement l’abréviation ensuite.

% Affiche la liste des acronymes
\printglossary[
  type=\acronymtype,
  title=Acronymes,
  toctitle=Acronymes % Alternative : Liste des acronymes, changera le titre uniquement dans la toc (table of content)
]

% Affiche le glossaire principal
\printglossary[
  title=Glossaire,
  toctitle=Glossaire % Glossaire des termes
]

% Utilisation [à supprimer]
Une fois un mot ou acronyme déclaré avec \newglossaryentry ou \newacronym :
Pour l’utiliser dans ton texte :
\gls{nmr} → affiche NMR + explication complète la 1re fois (si configuré)

\Gls{nmr} → idem, mais avec majuscule

\glspl{nmr} → pluriel

\acrshort{nmr} → affiche juste NMR

\acrfull{nmr} → affiche acide abscissique (ABA) (forme longue + sigle)

\acrlong{nmr} → juste la forme longue

LaTeX se charge automatiquement d’écrire le mot complet la première fois, puis seulement l’abréviation ensuite.

% Corps principal
\mainmatter
\addcontentsline{toc}{chapter}{Introduction générale}
\chapter*{Introduction générale}

Lorem ipsum dolor sit amet, consectetur adipiscing elit, sed do eiusmod tempor incididunt ut labore et dolore magna aliqua. Ut enim ad minim veniam, quis nostrud exercitation ullamco laboris nisi ut aliquip ex ea commodo consequat. Duis aute irure dolor in reprehenderit in voluptate velit esse cillum dolore eu fugiat nulla pariatur. Excepteur sint occaecat cupidatat non proident, sunt in culpa qui officia deserunt mollit anim id est laborum.

Lorem ipsum dolor sit amet, consectetur adipiscing elit, sed do eiusmod tempor incididunt ut labore et dolore magna aliqua. Ut enim ad minim veniam, quis nostrud exercitation ullamco laboris nisi ut aliquip ex ea commodo consequat. Duis aute irure dolor in reprehenderit in voluptate velit esse cillum dolore eu fugiat nulla pariatur. Excepteur sint occaecat cupidatat non proident, sunt in culpa qui officia deserunt mollit anim id est laborum.
\include{1_Chapitre_1_Etat_Art/chapitre_1_etat_art}
%\include{2_Chapitre_2_Mat_et_Met/chapitre_2_materiel_et_methode}
% ...
\include{0_Liminaires/conclusion}

%\backmatter
% Bibliographie
\newpage
\addcontentsline{toc}{chapter}{Bibliographie} % ajoute manuellement la Bibliographie dans la table des matières (utile si \printbibliography ne le fait pas automatiquement).
\printbibliography % imprime ta bibliographie complète.

% Liste des figures
\newpage
\addcontentsline{toc}{chapter}{Table des figures}
\listoffigures

% Leste des tableaux
\newpage
\addcontentsline{toc}{chapter}{Liste des tableaux}
\listoftables

% Annexes
\appendix
% annexes
\chapter{Annexe Compléments sur les calculs}
Contenu...

\chapter{Annexe Script Python}
Contenu...

\end{document}

\section{Thèse sur article, dite aussi "thèse sur travaux" :}

Ce type de thèse présente les principaux résultats de la recherche sous forme d’articles publiés dans des revues scientifiques, ou sous forme d’articles acceptés pour publication, dans le domaine de recherche du doctorant.

 • Le doctorant peut être l’auteur unique d’un article ou son auteur principal, ce qui signifie :
 - qu’il précise, dans le cas d’articles co-signés, sa contribution qui doit être significative,
 - que les articles correspondent à une contribution scientifique originale,

 • Les articles qui constituent le corps du document doivent être :
 - précédés dans le mémoire de thèse, d’une présentation substantielle,
 - suivis d’une synthèse substantielle qui comprend une discussion générale des résultats, des conclusions et une bibliographie.

Les parties nécessaires à la compréhension générale du mémoire mais qui ne peuvent pas paraître dans les journaux scientifiques doivent être fournies (techniques ou méthodes d’analyse spécifiques, données supplémentaires, etc.).

\section{Redaction du manuscrit :}
La rédaction du mémoire est l’une des formations apportées par le doctorat et constitue une part essentielle du travail de thèse.
Le mémoire doit montrer les capacités du candidat à présenter l’objectif de son étude, faire l’état de l’art sur la question, situer son travail dans le contexte international, exposer la méthodologie, développer les différents aspects de son travail. Ses capacités à rédiger seront utiles au docteur au cours de sa vie professionnelle, qu’elle se réalise dans le milieu académique ou en entreprise.

\subparagraph{Contenu :}

La forme traditionnelle du mémoire est constituée de différents chapitres que le doctorant s’attachera à articuler entre eux. Il présentera une synthèse critique des résultats obtenus avant de proposer des conclusions et des perspectives.Le mémoire pourra intégrer des articles publiés ou acceptés pour publication dans des revues reconnues dans le domaine de recherche du doctorant à condition :
• qu’il en soit l’artisan principal et qu’il précise, dans le cas d’articles co-signés, sa contribution qui doit être significative,
• que les articles correspondent à des contributions scientifiques originales,
• qu’ils soient précédés, dans le mémoire de thèse, d’une présentation substantielle montrant comment ils s’intègrent dans le travail de thèse,
• que les parties nécessaires à la compréhension générale du mémoire mais qui ne peuvent pas paraître dans les journaux scientifiques soient fournies (techniques ou méthodes d’analyse spécifiques, données supplémentaires, etc.).

Dans le cas où de larges passages du manuscrit seraient constitués d'articles, il est demandé un travail de rédaction d'au minimum 40 pages introduisant le sujet, présentant l'état de l'art et toute partie ne figurant pas dans les articles mais étant nécessaire à la compréhension de l’ensemble, liant les articles entre eux et concluant sur les objectifs atteints et les perspectives.
Des publications ne sauraient, en aucun cas, représenter à elles seules la thèse elle-même.

Le rôle des membres du jury sera précisé après soutenance dans la version définitive du manuscrit de thèse.

\subparagraph{L’usage des langues pour la rédaction du mémoire et la soutenance de thèse :}

La langue de rédaction est normalement le français.

Le mémoire contient obligatoirement un résumé rédigé en français et en anglais.

La langue de soutenance est normalement le français, sauf pour les étudiants non francophones, ou ceux dont le travail s'est développé dans un cadre international, auquel cas l'usage de la langue anglaise est toléré.
La discussion pourra se dérouler en anglais si des membres du jury le souhaitent.
Les rapporteurs et le président du jury devront rédiger leurs rapports en suivant ce qui est prévu dans la convention de thèse si elle existe, en français ou en anglais.

 